\documentclass[fleqn]{article}
\usepackage[utf8]{inputenc}
\usepackage[T2A]{fontenc}
\usepackage[russian]{babel}
\usepackage{graphicx}
\usepackage{amsmath}
\renewcommand{\arraystretch}{1.5}
\renewcommand{\baselinestretch}{1.75}
\usepackage[margin=3cm]{geometry}

\usepackage{amssymb}
\usepackage{caption}
\setlength{\mathindent}{0pt}

\usepackage{titlesec} % уменьшение отступа над главами
\titlespacing*{\section}{0pt}{0pt}{1em}  % уменьшение отступа над главами 



\begin{document}

\title{ВЛИЯНИЕ ПРОМЫСЛА НА СЦЕНАРИИ РАЗВИТИЯ ДВУХВОЗРАСТНОЙ ПОПУЛЯЦИИ ПРИ ПЛОТНОСТНОЙ РЕГУЛЯЦИИ РОЖДАЕМОСТИ}
\author{Я}
\maketitle

\tableofcontents
\newpage

\section{Постановка модели}

\[
\begin{cases}
X_{n+1} = \alpha Y_n(1-\beta Y_n)(1-u_1) \\
Y_{n+1} = (sX_n + vY_n)(1-u_2)
\end{cases}
\]
Найду неподвижные точки. То есть $X_{n+1} = X_n = X$ и $Y_{n+1} = Y_n = Y$
\[
\begin{cases}
X = \alpha Y(1-\beta Y)(1-u_1) \\
Y = (sX_n + vY)(1-u_2)
\end{cases}
\]
Подставляю X во второе уравнение \\
$Y = (s \alpha Y (1-\beta Y) (1-u_1) + v Y) (1-u_2) \\$
$ Y = ((s \alpha Y - s \alpha \beta Y^2)(1-u_1) + v Y)(1-u_2) | :Y( Y \neq 0) \\$
$ 1 = ((s\alpha - s\alpha \beta Y)(1-u_1) + v)(1-u_2)  \\$
$ 1 = (s \alpha - s\alpha\beta Y - s\alpha u_1 + s \alpha \beta u_1 Y + v)(1-u_2) \\$
$1 = s\alpha - s\alpha\beta Y - s\alpha u_1 + s\alpha\beta u_1 Y + v 
    - s\alpha u_2 + s\alpha\beta u_2 Y + s\alpha\beta u_1 u_2 
    - s\alpha u_1 u_2 Y - v u_2 \\$
$ Y = \dfrac{-1 + s \alpha  - s \alpha u_1 + v -s \alpha u_2 + s \alpha u_1 u_2 - v u_2}{s \alpha \beta - s \alpha \beta u_1 - s \alpha \beta u_2 + s \alpha \beta u_1 u_2} \\$
$ Y = \dfrac{s \alpha (1-u_1)(1-u_2) + v(1-u_2)-1}{s \alpha \beta (1-u_1)(1-u_2)} $ \\
\\
$Y = \dfrac{s \alpha (1-u_1)(1-u_2) + v(1-u_2)-1}{s \alpha \beta (1-u_1)(1-u_2)}$ \\ 
$X = \dfrac{\big(s\alpha(1-u_1)(1-u_2) + v(1-u_2) - 1\big)\big(1 - v(1-u_2)\big)} {s^2\alpha\beta(1-u_1)(1-u_2)^2} $ \\
Теперь выполню переход к новой системе координат: $ y = \beta Y; x = s \beta X; \tilde{\alpha} = \alpha s$. Получаю: \\
$ y = 1 + \dfrac{v(1-u_2) - 1}{ \tilde{\alpha}(1-u_1)(1-u_2)} $ \\
$ x = \dfrac{y(1-v(1-u_2))}{1-u_2} $ \\
В силу условие Ферхюста, накладываются ограничения на y: $ 0 \leq y\leq  1 $.\\
После перехода получил из 6 параметров 4. \\
\\

\section{Анализ устойчивости}

 Возьму изначальные равенства. Без n так как я пока что рассматриваю неподвижные точки, не зависящие от времени \\
$ F_1(X,Y) = \alpha Y_n(1-\beta Y_n)(1-u_1) \\ $
$ F_1(X,Y) = \alpha (Y_n-\beta Y_n^2)(1-u_1) \\ $
$ F_2(X,Y) = (sX_n + vY_n)(1-u_2) $ \\
Построю матрицу Якоби или же матрицу возмущений:
 $$ J = \begin{pmatrix} 
    F'_{1X} & F'_{1Y} \\ 
    F'_{2X} & F'_{2Y}
 \end{pmatrix}
  = \begin{pmatrix} 
    0 & \alpha(1-u_1)(1-2\beta Y_n) \\ 
    s(1-u_2) & v(1-u_2)
 \end{pmatrix} \overset{\text{Переобозначил}}{=} 
 \begin{pmatrix} 
    0 & b \\ 
    c & d
 \end{pmatrix}
 \\$$

 И при этом собственные числа это мультипликаторы отображения $\lambda_1$ и $ \lambda_2$, для которых справедливо: $ \lambda^2 - trJ \lambda + |J| = 0 $ \\
 $ p = trJ = d  $ (след)\\
 $ q = |J| = -bc $\\
 Характеристическое уравнение: $ \lambda^2 - p \lambda + q = 0 $\\
 Характеристическое уравнение: $ \lambda^2 - d \lambda - bc = 0 $\\
 $ D = d^2 + 4bc $
 Теперь построю треугольник устойчивости, краями которого есть 3 прямые: \\
   1) $ q = p - 1 \quad (\lambda = 1)$ - \textbf{транскритическая бифуркация}\\
   2) $ q = -1 - p \quad (\lambda = -1)$ - \textbf{бифуркация удвоения периода (flip)}\\
   3) $ q = 1 \quad $ ($\lambda$ - комплексное) - \textbf{бифуркация Неймарка–Сакера} \\
Решение будет устойчиво при всех  $|\lambda| < 1$, то есть устойчивые решения находятся внутри этого треугольника \\
\\
$\tilde{\alpha} = s \alpha$ --- сочетание рождаемости и коэффициента перехода между возрастами.\\
$0 \leq v \leq 1$ --- выживаемость взрослых.\\
$0 \leq u_1 \leq 1$ --- доля вылова малых рыб.\\
$0 \leq u_2 \leq 1$ --- доля вылова взрослых рыб.\\
1. Рассмотрю прямую $q = p - 1$ \\
$ - \alpha(1-u_2) \alpha (1-u_1) (1-2\beta Y_n)  = v(1-u_2)- 1$ \\
$ - \tilde{\alpha} (1-u_2)(1-u_1)(1-2y) = v(1-u_2) - 1$ \\
Теперь зафиксирую 


\end{document}